% Options for packages loaded elsewhere
\PassOptionsToPackage{unicode}{hyperref}
\PassOptionsToPackage{hyphens}{url}
\PassOptionsToPackage{dvipsnames,svgnames*,x11names*}{xcolor}
%
\documentclass[
  11pt,
]{article}
\usepackage{lmodern}
\usepackage{amssymb,amsmath}
\usepackage{ifxetex,ifluatex}
\ifnum 0\ifxetex 1\fi\ifluatex 1\fi=0 % if pdftex
  \usepackage[T1]{fontenc}
  \usepackage[utf8]{inputenc}
  \usepackage{textcomp} % provide euro and other symbols
\else % if luatex or xetex
  \usepackage{unicode-math}
  \defaultfontfeatures{Scale=MatchLowercase}
  \defaultfontfeatures[\rmfamily]{Ligatures=TeX,Scale=1}
\fi
% Use upquote if available, for straight quotes in verbatim environments
\IfFileExists{upquote.sty}{\usepackage{upquote}}{}
\IfFileExists{microtype.sty}{% use microtype if available
  \usepackage[]{microtype}
  \UseMicrotypeSet[protrusion]{basicmath} % disable protrusion for tt fonts
}{}
\makeatletter
\@ifundefined{KOMAClassName}{% if non-KOMA class
  \IfFileExists{parskip.sty}{%
    \usepackage{parskip}
  }{% else
    \setlength{\parindent}{0pt}
    \setlength{\parskip}{6pt plus 2pt minus 1pt}}
}{% if KOMA class
  \KOMAoptions{parskip=half}}
\makeatother
\usepackage{xcolor}
\IfFileExists{xurl.sty}{\usepackage{xurl}}{} % add URL line breaks if available
\IfFileExists{bookmark.sty}{\usepackage{bookmark}}{\usepackage{hyperref}}
\hypersetup{
  colorlinks=true,
  linkcolor=blue,
  filecolor=Maroon,
  citecolor=Blue,
  urlcolor=Blue,
  pdfcreator={LaTeX via pandoc}}
\urlstyle{same} % disable monospaced font for URLs
\usepackage[margin=1in]{geometry}
\usepackage{longtable,booktabs}
% Correct order of tables after \paragraph or \subparagraph
\usepackage{etoolbox}
\makeatletter
\patchcmd\longtable{\par}{\if@noskipsec\mbox{}\fi\par}{}{}
\makeatother
% Allow footnotes in longtable head/foot
\IfFileExists{footnotehyper.sty}{\usepackage{footnotehyper}}{\usepackage{footnote}}
\makesavenoteenv{longtable}
\setlength{\emergencystretch}{3em} % prevent overfull lines
\providecommand{\tightlist}{%
  \setlength{\itemsep}{0pt}\setlength{\parskip}{0pt}}
\setcounter{secnumdepth}{-\maxdimen} % remove section numbering

\author{Group AN — North America + Europe}
\date{May 2026}

\begin{document}
\maketitle
\tableofcontents
\newpage




\hypertarget{data-loading-and-investment-sets}{%
\section{1. Data Loading and Investment
Sets}\label{data-loading-and-investment-sets}}

\hypertarget{data-sources}{%
\subsection{1.1 Data Sources}\label{data-sources}}

The project uses processed investment set files (CSV format) for each
year from 2013 to 2025. Each file, named investment\_set\_\{year\}.csv,
contains one row per firm and the following key variables:

\begin{longtable}[]{@{}ll@{}}
\toprule
\textbf{Variable} & \textbf{Description}\tabularnewline
\midrule
\endhead
ISIN & Firm identifier\tabularnewline
NAME, Country, Region & Firm identifiers and geographic
classification\tabularnewline
MV\_Y & End-of-year market capitalisation (million USD)\tabularnewline
CO2\_S1, CO2\_S2 & Scope 1 and Scope 2 CO₂ emissions
(tonnes)\tabularnewline
REV & Annual revenues (thousands USD)\tabularnewline
CI & Carbon intensity = (CO2\_S1+CO2\_S2)/(REV/1000) (tCO₂e per million
USD revenue)\tabularnewline
Monthly return columns & 120 monthly total returns (10-year estimation
window ending in year Y)\tabularnewline
\bottomrule
\end{longtable}

Carbon data covers 2002--2024. Due to limited coverage before 2010, the
allocation exercise begins at end of 2013 (portfolio first invested in
2014).

\hypertarget{investment-set-construction}{%
\subsection{1.2 Investment Set
Construction}\label{investment-set-construction}}

For each year Y = 2013, \ldots, 2024, the investment set is defined as
firms satisfying all of the following criteria:

\begin{itemize}
\item
  Belong to the assigned region.
\item
  Have sufficient return history: at least 3 years (36 months) of
  non-missing monthly returns over the 10-year window ending at end of
  year Y.
\item
  Have carbon data (CO₂ and revenues) available at end of year Y.
\item
  Do not have stale prices: firms whose monthly returns are zero for
  more than 50\% of the estimation window are excluded.
\end{itemize}

The investment decision is made at end of year Y for portfolio
implementation over year Y+1. The covariance matrix and expected returns
are estimated from the 10-year window of monthly returns available at
end of Y.

\hypertarget{full-available-data-variant}{%
\subsection{1.3 Full-Available-Data
Variant}\label{full-available-data-variant}}

A second version of each investment set, labelled is\_\{year\}\_fa, is
constructed by dropping any firms with at least one missing monthly
return over the 10-year window. This stricter filter is used for the
full-available-data value-weighted portfolio and as the benchmark in the
Ledoit-Wolf tracking-error and net-zero optimizations.

\hypertarget{value-weighted-portfolios}{%
\section{2. Value-Weighted Portfolios}\label{value-weighted-portfolios}}

\hypertarget{standard-value-weighted-portfolio}{%
\subsection{2.1 Standard Value-Weighted
Portfolio}\label{standard-value-weighted-portfolio}}

At the end of each year Y, portfolio weights are set proportional to
each firm's market capitalisation within the investment set:

\emph{w\_\{i,Y\} = MV\_\{i,Y\} / Σ\_j MV\_\{j,Y\}}

These end-of-year weights are stored in the column MV\_Y\_weight and
serve as the January weight for year Y+1. However, rather than holding
these weights fixed throughout the year, the portfolio is reweighted
each month using the actual end-of-month market capitalisations from
MV\_M\_filtered.csv. Concretely, for February through December of year
Y+1, the weight of firm i at the start of month t is:

\emph{w\_\{i,t\} = MV\_\{i,t-1\} / Σ\_j MV\_\{j,t-1\} for t = Feb(Y+1),
\ldots, Dec(Y+1)}

where MV\_\{i,t-1\} is the market capitalisation of firm i at the end of
the preceding month, taken from MV\_M\_filtered.csv. Monthly caps are
matched to return columns by nearest available date. If a monthly cap is
unavailable, the end-of-year market cap (MV\_Y) from the prior
investment set is used as a fallback. Weights are renormalized within
the investment-set universe so they sum to 1 each month. The monthly
portfolio return is then R\_\{p,t\} = Σ\_i w\_\{i,t\} × R\_\{i,t\},
where missing returns are replaced by zero to avoid survivorship bias
(firms that delist or have gaps are treated as contributing 0 return
rather than being dropped from the denominator).

\hypertarget{full-available-data-value-weighted-portfolio}{%
\subsection{2.2 Full-Available-Data Value-Weighted
Portfolio}\label{full-available-data-value-weighted-portfolio}}

The same monthly market cap reweighting methodology is applied to the
is\_\{year\}\_fa subset (firms with complete 10-year return histories).
After removing firms with any NaN in their return series, the
end-of-year MV\_Y\_weight values no longer sum to 1 over the reduced
universe; they are therefore explicitly renormalized so that Σ\_i
w\_\{i,Y\}\^{}\{fa\} = 1. Monthly reweighting from February onwards then
uses the same MV\_M\_filtered.csv caps restricted to the fa universe,
renormalized within that universe. This portfolio is stored as
MV\_Y\_weight\_fa and serves as the benchmark for tracking-error and
net-zero optimizations.

\hypertarget{minimum-variance-portfolio-optimization}{%
\section{3. Minimum-Variance Portfolio
Optimization}\label{minimum-variance-portfolio-optimization}}

\hypertarget{covariance-matrix-estimation}{%
\subsection{3.1 Covariance Matrix
Estimation}\label{covariance-matrix-estimation}}

For each year Y, two covariance matrices are computed from the 10-year
window of monthly returns:

Sample Covariance (SLSQP method): The standard sample covariance matrix
Σ\_Y = (1/τ) Σ\_\{k=0\}\^{}\{τ-1\} (R\_\{t-k\} − μ\_Y)(R\_\{t-k\} −
μ\_Y)' where τ = 120 months. Computed using pandas .cov() on the
transposed return matrix.

Ledoit-Wolf Shrinkage (CVXPY method): Before applying Ledoit-Wolf, firms
with any NaN in their return series are removed. The remaining returns
are passed to scikit-learn's LedoitWolf estimator, which produces a
positive-semidefinite covariance matrix Σ\_Y\^{}\{LW\} by optimally
shrinking the sample covariance towards a scaled identity. The resulting
subset of firms is stored as returns\_\{year\}\_lw.

\hypertarget{slsqp-optimization-sample-covariance}{%
\subsection{3.2 SLSQP Optimization (Sample
Covariance)}\label{slsqp-optimization-sample-covariance}}

The minimum-variance weights for year Y are obtained by solving:

\emph{min\_\{α\} α' Σ\_Y α}

\emph{s.t. Σ\_i α\_i = 1, α\_i ≥ 0 for all i}

Solved with scipy.optimize.minimize using the SLSQP method. A numerical
scaling factor of 10,000 is applied to Σ\_Y before optimization to
improve convergence, and the resulting weights are verified to be
identical to those from the unscaled problem. The initial guess is the
equal-weight vector.

\hypertarget{cvxpy-osqp-optimization-ledoit-wolf}{%
\subsection{3.3 CVXPY / OSQP Optimization
(Ledoit-Wolf)}\label{cvxpy-osqp-optimization-ledoit-wolf}}

The same quadratic program is solved with the CVXPY library using the
OSQP solver on the Ledoit-Wolf covariance matrix Σ\_Y\^{}\{LW\}:

\emph{min\_w w' Σ\_Y\^{}\{LW\} w (CVXPY quad\_form)}

\emph{s.t. Σ\_i w\_i = 1, w\_i ≥ 0}

A scaling factor of 10,000 is also applied for numerical stability.
Small negative weights produced by solver numerical noise are clipped to
zero and weights are renormalized to sum to 1.

\hypertarget{out-of-sample-returns}{%
\subsection{3.4 Out-of-Sample Returns}\label{out-of-sample-returns}}

For each year Y = 2013, \ldots, 2024, the weights α\_Y computed at end
of Y are applied to the actual monthly returns of year Y+1. The
portfolio is rebalanced annually; within the year, weights drift with
prices (dynamic weight update). This procedure produces a 144-month
out-of-sample return time series from January 2014 to December 2025.

\hypertarget{carbon-metrics}{%
\section{4. Carbon Metrics}\label{carbon-metrics}}

\hypertarget{carbon-intensity-ci}{%
\subsection{4.1 Carbon Intensity (CI)}\label{carbon-intensity-ci}}

Carbon intensity measures CO₂ emissions relative to revenues. For each
firm i in year Y:

\emph{CI\_\{i,Y\} = (CO2\_S1\_\{i,Y\} + CO2\_S2\_\{i,Y\}) /
(REV\_\{i,Y\} / 1000)}

The division by 1000 converts revenues from thousands to millions of
USD, giving CI in tonnes of CO₂ equivalent per million USD of revenue
(tCO₂e / M\$rev).

\hypertarget{weighted-average-carbon-intensity-waci}{%
\subsection{4.2 Weighted Average Carbon Intensity
(WACI)}\label{weighted-average-carbon-intensity-waci}}

The portfolio-level WACI aggregates firm-level carbon intensity using
portfolio weights:

\emph{WACI\_Y\^{}\{(p)\} = Σ\_i α\_\{i,Y\} × CI\_\{i,Y\}}

Computed annually for each portfolio using the weights determined at end
of year Y.

\hypertarget{carbon-footprint-cf}{%
\subsection{4.3 Carbon Footprint (CF)}\label{carbon-footprint-cf}}

The carbon footprint of the portfolio measures the CO₂ emissions
attributed to the investor per million USD invested:

\emph{CF\_Y\^{}\{(p)\} = Σ\_i α\_\{i,Y\} × E\_\{i,Y\} / Cap\_\{i,Y\}}

where E\_\{i,Y\} = CO2\_S1\_\{i,Y\} + CO2\_S2\_\{i,Y\} (total emissions
in tonnes), Cap\_\{i,Y\} = MV\_Y (market capitalisation in million USD),
and α\_\{i,Y\} is the portfolio weight. The result is in tonnes of CO₂
per million USD invested (tCO₂e / M\$inv). Portfolio wealth cancels out
of this formula, so no wealth evolution is required.

\hypertarget{cf-constrained-minimum-variance-portfolio-part-3.2}{%
\section{5. CF-Constrained Minimum-Variance Portfolio (Part
3.2)}\label{cf-constrained-minimum-variance-portfolio-part-3.2}}

A carbon footprint constraint is added to the minimum-variance problem
to force the portfolio CF to be no more than 50\% of the unconstrained
minimum-variance portfolio CF in the same year:

\emph{min\_\{α\} α' Σ\_Y α}

\emph{s.t. CF\_Y\^{}\{(p)\} ≤ 0.5 × CF\_Y\^{}\{(mv)\}}

\emph{Σ\_i α\_i = 1, α\_i ≥ 0}

where CF\_Y\^{}\{(mv)\} is the carbon footprint of the unconstrained
minimum-variance portfolio computed for year Y. The constraint is
re-evaluated each year so that the 50\% reduction target is relative to
the current unconstrained portfolio, not a fixed baseline.

Two implementations are produced:

\begin{itemize}
\item
  SLSQP (sample covariance): same solver as Section 3.2, warm-started
  from equal weights.
\item
  CVXPY/OSQP (Ledoit-Wolf): same solver as Section 3.3, with the CF
  constraint expressed as a linear inequality e' w ≤ target where e\_i =
  E\_\{i,Y\}/Cap\_\{i,Y\}.
\end{itemize}

\hypertarget{tracking-error-minimization-part-3.3}{%
\section{6. Tracking Error Minimization (Part
3.3)}\label{tracking-error-minimization-part-3.3}}

This strategy minimizes the tracking error relative to the
value-weighted benchmark while constraining the portfolio carbon
footprint to 50\% of the VW benchmark CF. It targets a passive investor
seeking to reduce carbon emissions with minimal deviation from the
market-cap portfolio.

The optimization problem is:

\emph{min\_\{α\} TE\_Y\^{}2 = (α\_Y − α\_Y\^{}\{vw\})' Σ\_Y (α\_Y −
α\_Y\^{}\{vw\})}

\emph{s.t. CF\_Y\^{}\{(p)\} ≤ 0.5 × CF\_Y\^{}\{(vw)\}}

\emph{Σ\_i α\_i = 1, α\_i ≥ 0}

Minimizing TE² is equivalent to minimizing TE since the square-root
function is strictly increasing over non-negative values and is
numerically preferable near zero.

Two implementations:

\begin{itemize}
\item
  SLSQP: uses the sample covariance matrix and MV\_Y\_weight as the VW
  benchmark. Warm-started from the VW weights to ensure feasibility.
\item
  CVXPY/OSQP with Ledoit-Wolf: uses Σ\_Y\^{}\{LW\} and MV\_Y\_weight\_fa
  (the full-available-data VW benchmark). The VW\_fa weights are
  reindexed to the LW subset of firms and renormalized to sum to 1
  before use in the objective. The CF reference is computed over the
  full investment set using MV\_Y\_weight\_fa.
\end{itemize}

\hypertarget{net-zero-portfolio-part-4.1}{%
\section{7. Net Zero Portfolio (Part
4.1)}\label{net-zero-portfolio-part-4.1}}

The net zero strategy adopts the same tracking-error minimization
framework as Section 6, but replaces the static 50\% CF reduction
constraint with a time-varying constraint that enforces a compounding
10\% annual reduction in the portfolio carbon footprint from a fixed
2013 baseline:

\emph{CF\_Y\^{}\{(p)\} ≤ (1 − θ)\^{}\{Y−Y\_0+1\} ×
CF\_\{Y\_0\}\^{}\{(vw)\}}

where θ = 0.10, Y\_0 = 2013, and CF\_\{2013\}\^{}\{(vw)\} is the carbon
footprint of the value-weighted portfolio in 2013 (computed once as a
fixed reference). This gives the following annual carbon budgets:

\begin{longtable}[]{@{}ll@{}}
\toprule
\textbf{Year Y} & \textbf{CF budget as \% of 2013 VW CF}\tabularnewline
\midrule
\endhead
2013 (Y−Y₀+1 = 1) & 90.0\% (0.9¹)\tabularnewline
2014 (exponent = 2) & 81.0\% (0.9²)\tabularnewline
2016 (exponent = 4) & 65.6\% (0.9⁴)\tabularnewline
2019 (exponent = 7) & 47.8\% (0.9⁷)\tabularnewline
2022 (exponent = 10) & 34.9\% (0.9¹⁰)\tabularnewline
2024 (exponent = 12) & 28.2\% (0.9¹²)\tabularnewline
\bottomrule
\end{longtable}

The net-zero strategy is implemented using the CVXPY/Ledoit-Wolf solver
only. The SLSQP implementation was removed as the LW formulation
provides a more robust and numerically stable solution for this
constrained problem.

\begin{itemize}
\item
  CVXPY/LW: CF reference = 2013 VW\_fa CF using MV\_Y\_weight\_fa
  weights computed over the fa universe. CVXPY weights are stored with
  .flatten() to ensure a 1-D array regardless of solver output shape.
\item
  The nz\_weight\_slsqp column is no longer populated in the investment
  set CSVs; only nz\_weight\_lw is stored.
\end{itemize}

CVXPY weights are stored with .flatten() to ensure a 1-D array
regardless of solver output shape, avoiding downstream dtype errors.

\hypertarget{portfolio-weight-storage-and-csv-management}{%
\section{8. Portfolio Weight Storage and CSV
Management}\label{portfolio-weight-storage-and-csv-management}}

All portfolio weights are collected into a unified per-year investment
set DataFrame (inv\_set\_\{year\}) and persisted to CSV to allow
inter-cell continuity without re-running the full optimization pipeline.
The final column structure is:

\begin{longtable}[]{@{}ll@{}}
\toprule
\textbf{Column Index} & \textbf{Contents}\tabularnewline
\midrule
\endhead
0--4 & ISIN, NAME, Country, Region, MV\_Y\tabularnewline
5 & MV\_Y\_weight (value-weighted)\tabularnewline
6 & MV\_Y\_weight\_fa (full-available-data VW)\tabularnewline
7 & min\_var\_weight\_slsqp (unconstrained MV, SLSQP)\tabularnewline
8 & min\_var\_weight\_lw (unconstrained MV, LW)\tabularnewline
9 & min\_var\_weight\_cf\_slsqp (CF-constrained MV,
SLSQP)\tabularnewline
10 & min\_var\_weight\_cf\_lw (CF-constrained MV, LW)\tabularnewline
11 & te\_weight\_slsqp (tracking error, SLSQP)\tabularnewline
12 & te\_weight\_lw (tracking error, LW)\tabularnewline
13 & nz\_weight\_slsqp (net zero, SLSQP --- removed, column not
populated)\tabularnewline
14 & nz\_weight\_lw (net zero, LW)\tabularnewline
15+ & CO2\_S1, CO2\_S2, REV, CI, and monthly return
columns\tabularnewline
\bottomrule
\end{longtable}

When adding LW-specific weights (which are only defined for the NaN-free
LW subset), a left-merge on ISIN is performed so that firms outside the
LW universe receive a weight of 0. Stale columns from previous runs are
dropped before merging to avoid pandas 2.0+ MergeError on duplicate
column names.

\hypertarget{out-of-sample-return-computation}{%
\section{9. Out-of-Sample Return
Computation}\label{out-of-sample-return-computation}}

\hypertarget{return-dataframe-format}{%
\subsection{9.1 Return Dataframe Format}\label{return-dataframe-format}}

Returns are stored in a wide format using the create\_returns\_dataframe
helper: a DataFrame with years (2014--2025) as the row index and months
(Jan--Dec) as the twelve columns. A monthly time series is obtained by
calling .stack() on this DataFrame, producing a Series indexed by (year,
month) tuples.

\hypertarget{value-weighted-portfolio-returns}{%
\subsection{9.2 Value-Weighted Portfolio
Returns}\label{value-weighted-portfolio-returns}}

For the value-weighted portfolios (both full investment set and fa
subset), the return for each month of year Y+1 is computed using the
appropriate monthly weight column. Return columns and weight columns are
paired by sorting both independently by date, guaranteeing correct
alignment regardless of column order in the DataFrame. The steps are:

\begin{itemize}
\item
  Build a per-year DataFrame containing MV\_Y\_weight (January weight)
  and MV\_M\_weight\_\{date\} columns (February through November
  weights) from actual monthly market caps, plus the 12 monthly return
  columns from is\_\{year\}.
\item
  For each month t, multiply the corresponding weight column by the
  return column (with missing returns filled as zero) and sum across all
  firms to obtain R\_\{p,t\}.
\item
  December uses the November end-of-month cap weights (the 11th MV\_M
  column), which together with MV\_Y\_weight gives exactly 12
  weight-return pairings per year.
\end{itemize}

\hypertarget{minimum-variance-cf-constrained-tracking-error-and-net-zero-portfolio-returns}{%
\subsection{9.3 Minimum-Variance, CF-Constrained, Tracking-Error, and
Net-Zero Portfolio
Returns}\label{minimum-variance-cf-constrained-tracking-error-and-net-zero-portfolio-returns}}

All optimized portfolios (SLSQP, LW, CF-constrained, tracking-error, and
net-zero) use a buy-and-hold weight drift approach rather than fixed
monthly weights. Weights are set once at end of year Y and then allowed
to evolve naturally as asset prices move, exactly as they would for an
investor who invests at the start of the year and does not rebalance
intra-year. The update rule for each month t within year Y+1 is:

\begin{itemize}
\item
  w\_\{i,t\} = w\_\{i,t-1\} × (1 + R\_\{i,t-1\}) / (1 + R\_\{p,t-1\})
\item
  An asset whose return exceeds the portfolio return this month will
  have a higher weight next month; an asset that underperforms will
  shrink in share. The division by (1 + R\_\{p,t-1\}) ensures that the
  updated weights still sum to 1 at all times. Missing returns are
  filled as zero so that delisted or temporarily absent firms contribute
  0 to portfolio return without distorting the weight denominator (no
  survivorship bias).
\item
  This procedure is applied identically to all optimized strategies:
  SLSQP min-variance, LW min-variance, CF-constrained (SLSQP and LW),
  tracking-error (SLSQP and LW), and net-zero (LW only). The result is a
  12-element monthly return series per year, stacked across years to
  form the full 144-month out-of-sample series (January 2014 to December
  2025).
\end{itemize}

\hypertarget{performance-statistics}{%
\subsection{9.4 Performance Statistics}\label{performance-statistics}}

For each portfolio, the following statistics are computed over the full
144-month period (January 2014 -- December 2025):

\begin{longtable}[]{@{}ll@{}}
\toprule
\textbf{Statistic} & \textbf{Formula}\tabularnewline
\midrule
\endhead
Annualized Average Return & (1 + mean(r\_monthly))\^{}12 −
1\tabularnewline
Annualized Volatility & std(r\_monthly) × √12\tabularnewline
Sharpe Ratio & Annualized Return / Annualized Volatility\tabularnewline
Annualized Cumulative Return & (1 + cum\_return)\^{}\{1/12\} −
1\tabularnewline
Min/Max Monthly Return & min and max of r\_monthly\tabularnewline
\bottomrule
\end{longtable}

\hypertarget{carbon-footprint-and-waci-evolution}{%
\section{10. Carbon Footprint and WACI
Evolution}\label{carbon-footprint-and-waci-evolution}}

For each portfolio and each year Y = 2013, \ldots, 2024, the annual CF
and WACI are computed from the weights determined at end of Y and the
carbon data for year Y. For LW portfolios, the weight vector is
reindexed from the LW-subset order to the full inv\_set order using a
ISIN-based reindex, and emission intensities are computed over the
matching LW subset.

Results are collected into cf\_waci\_*\_df DataFrames indexed by year,
enabling line plots that show the evolution of both metrics across the
sample period for all portfolio variants.

\hypertarget{top-10-carbon-intensive-firms-analysis}{%
\section{11. Top-10 Carbon Intensive Firms
Analysis}\label{top-10-carbon-intensive-firms-analysis}}

For each year and each portfolio, the 10 firms with the highest carbon
intensity (CI) are identified from the investment set. Their portfolio
weights are displayed as horizontal bar charts in a 4×3 grid (one
subplot per year, 2013--2024). Each portfolio family uses a consistent
colour scheme:

\begin{longtable}[]{@{}ll@{}}
\toprule
\textbf{Portfolio Family} & \textbf{Bar Colour}\tabularnewline
\midrule
\endhead
Unconstrained MV (SLSQP and LW) & Steelblue\tabularnewline
CF-Constrained MV (SLSQP and LW) & Tomato\tabularnewline
Tracking Error (SLSQP and LW) & Medium Sea Green\tabularnewline
Net Zero (LW only) & Dark Orchid\tabularnewline
\bottomrule
\end{longtable}

These plots allow visual inspection of how much weight each strategy
allocates to the most carbon-intensive names, and how this evolves as
constraints tighten.

\hypertarget{decomposition-of-cf-and-waci-changes}{%
\section{12. Decomposition of CF and WACI
Changes}\label{decomposition-of-cf-and-waci-changes}}

A counterfactual decomposition separates the observed change in CF and
WACI from the 2013 level into two effects, following a Laspeyres-type
framework:

\hypertarget{emission-effect-frozen-weights}{%
\subsection{12.1 Emission Effect (frozen
weights)}\label{emission-effect-frozen-weights}}

The portfolio weights are held fixed at their 2013 values while
firm-level emissions and carbon intensity are allowed to evolve. This
isolates the contribution of firm-level decarbonization (or
deterioration) to the portfolio-level carbon metric, independently of
reweighting decisions.

\emph{CF\_Y\^{}\{emission\} = Σ\_i (w\_\{i,2013\} / Σ\_j w\_\{j,2013\})
× E\_\{i,Y\} / Cap\_\{i,Y\}}

\hypertarget{composition-effect-frozen-emissions}{%
\subsection{12.2 Composition Effect (frozen
emissions)}\label{composition-effect-frozen-emissions}}

The firm-level emissions and carbon intensities are held fixed at 2013
values while the portfolio weights are allowed to vary. This isolates
the contribution of active portfolio reweighting (tilting towards
lower-carbon firms) to the portfolio-level carbon metric.

\emph{CF\_Y\^{}\{composition\} = Σ\_i w\_\{i,Y\} × E\_\{i,2013\} /
Cap\_\{i,2013\}}

The same framework is applied to WACI, substituting CI for E/Cap. Each
decomposition figure contains two panels (WACI left, CF right), and is
produced for each portfolio variant. Only firms present in both the 2013
and year-Y investment sets (inner merge on ISIN) contribute to the
decomposition.

\hypertarget{summary-of-portfolio-families}{%
\section{13. Summary of Portfolio
Families}\label{summary-of-portfolio-families}}

\begin{longtable}[]{@{}ll@{}}
\toprule
\textbf{Portfolio} & \textbf{Description}\tabularnewline
\midrule
\endhead
VW & Value-weighted, standard MV\_Y weights\tabularnewline
VW (full available data) & Value-weighted, NaN-free firms only,
renormalized\tabularnewline
MV (SLSQP) & Min-variance, sample covariance, long-only\tabularnewline
MV (LW) & Min-variance, Ledoit-Wolf shrinkage, long-only\tabularnewline
MV CF-Constr. (SLSQP) & Min-variance + CF ≤ 50\% of MV CF (sample
cov.)\tabularnewline
MV CF-Constr. (LW) & Min-variance + CF ≤ 50\% of MV CF
(LW)\tabularnewline
TE (SLSQP) & Tracking-error min. + CF ≤ 50\% of VW CF (sample
cov.)\tabularnewline
TE (LW) & Tracking-error min. + CF ≤ 50\% of VW\_fa CF
(LW)\tabularnewline
NZ (LW only) & TE min. + CF ≤ (0.9)\^{}\{Y−2013+1\} ×
CF\_\{2013\}\^{}\{VW\_fa\} (LW, CVXPY/OSQP)\tabularnewline
NZ --- removed & SLSQP net-zero implementation removed; LW-only net-zero
listed above\tabularnewline
\bottomrule
\end{longtable}

Results

\textbf{Why the drastic decrease in CF and WACI?}

The issue is not a computation error but rather what the investment set
actually contains in each year. The investment set for 2013 is
restricted to firms that had carbon data available by end-2013 ---
mostly large, established companies in energy, materials, and
industrials, which are precisely the most carbon-intensive sectors. In
subsequent years, as carbon reporting improves and more firms start
disclosing emissions, many new companies enter the investment set. These
are predominantly tech, healthcare, and services companies --- sectors
with very low carbon intensity and, crucially, very high market
capitalizations.

\end{document}
